
\documentclass[11pt]{article}
\usepackage{amsmath, amsthm, amssymb}
\usepackage{quantikz}
\newtheorem{theorem}{Theorem}
\newtheorem{corollary}{Corollary}

% Success probability
\newcommand{\psucc}{P_{\textnormal{succ}}}
\newcommand{\perr}{P_{\textnormal{err}}}

% POVMs
\newcommand{\povm}[1]{\textnormal{POVM}_{\textnormal{#1}}}
\newcommand{\PiSet}[1]{\left\{ \Pi_1, \Pi_2, \dots, \Pi_{#1} \right\}}
\newcommand{\PiSetPrime}[1]{\left\{ \Pi_1', \Pi_2', \dots, \Pi_{#1}' \right\}}

% Superscripts
\newcommand{\supmed}{^{\textnormal{MED}}}
\newcommand{\supmedplus}{^{\textnormal{MED+}}}
\newcommand{\supuqsd}{^{\textnormal{OptUQSD}}}


\begin{document}

\begin{theorem} \label{thm:med_medplus}
    In the quantum state discrimination problem, given the same set of
    mixed quantum states $\{ \rho_{1}, \rho_{2}, \dots, \rho_{M} \}$
    and the same prior probabilities $\{p_{1}, p_{2}, \dots, p_{M}\}$, the probability of correctly identifying the given state in the optimal MED scenario is identical to that in the optimal MED+ scenario; i.e.,
    $(\psucc)\supmed = (\psucc)\supmedplus$.
    Furthermore, $\Pi_{M+1} = 0$ when the MED+ scenario becomes optimal.
\end{theorem}

\begin{proof}
    % Line break
    \noindent

    Let $\povm{M}$ be the collection of POVMs that each POVM has $M$ POVM
    elements, and here we allow a POVM element to be zero.
    Let $(\psucc)^{P}$ denote the probability of successful discrimination of
    some POVM $P$ or the POVM derived from some method $P$.
    We know that $(\psucc)\supmed >= (\psucc)^{P}, \; \forall \; P \in \povm{M}$.
    \begin{enumerate}
        \item $(\psucc)\supmed \leq (\psucc)\supmedplus$\\
              Assume $\povm{MED}$ is an optimal solution for MED, and\\
              $\povm{MED} \equiv \PiSet{M} \in \povm{M}$.
              We can construct $P \equiv \PiSetPrime{M+1} \in \povm{M+1}$
              by
              \begin{align*}
                  \Pi_{1}' = \Pi_{1}, \; \Pi_{2}' = \Pi_{2}, \; \dots,
                  \; \Pi_{M}' = \Pi_{M}, \; \Pi_{M+1}' = 0
              \end{align*}
              Then,
              \begin{align}
                  (\psucc)\supmedplus & \geq (\psucc)^{P}
                  = \sum_{i=1}^{M} p_{i} \operatorname{tr}(\rho_{i} \Pi_{i}')                      \\
                                      & = \sum_{i=1}^{M} p_{i} \operatorname{tr}(\rho_{i} \Pi_{i}) \\
                                      & = (\psucc)\supmed
              \end{align}
        \item $(\psucc)\supmed \geq (\psucc)\supmedplus$\\
              Assume $\povm{MED+}$ is an optimal solution for MED+, and\\
              $\povm{MED+} \equiv \PiSet{M+1} \in \povm{M+1}$.
              We can construct $P \equiv \PiSetPrime{M} \in \povm{M}$ by
              \begin{align*}
                  \Pi_{1}' = \Pi_{1}, \; \Pi_{2}' = \Pi_{2}, \; \dots,
                  \; \Pi_{M}' = \Pi_{M} + \Pi_{M+1}
              \end{align*}
              Then,
              \begin{align}
                  (\psucc)\supmed & \geq (\psucc)^{P}
                  = \sum_{i=1}^{M} p_{i} \operatorname{tr}(\rho_{i} \Pi_{i}')                           \\
                                  & = \sum_{i=1}^{M} p_{i} \operatorname{tr}(\rho_{i} \Pi_{i})
                  + p_{M} \operatorname{tr}(\rho_{i} \Pi_{M+1})                                         \\
                                  & = (\psucc)\supmedplus + p_{M} \operatorname{tr}(\rho_{i} \Pi_{M+1}) \\
                                  & \geq (\psucc)\supmedplus \label{eq:geq}
              \end{align}
        \item By 1. and 2., we can see that $(\psucc)\supmed = (\psucc)\supmedplus$.
        \item To make the equal sign in Eq (\ref{eq:geq}) hold, $\Pi_{M+1} = 0$.
    \end{enumerate}
\end{proof}

From Theorem \ref{thm:med_medplus}, we can easily prove the
following corollary.
\begin{corollary}
    In the quantum state discrimination problem, given the same set of
    mixed quantum states $\{ \rho_{1}, \rho_{2}, \dots, \rho_{M} \}$
    and the same prior probabilities $\{p_{1}, p_{2}, \dots, p_{M}\}$, the probability of correctly identifying the given state in the optimal MED scenario is always equal to or larger than that in the Optimal UQSD scenario; i.e.,
    $(\psucc)\supmed \geq (\psucc)\supuqsd$.
\end{corollary}

\begin{proof}
    If we see the MED+ problem and the Optimal UQSD problem as an convex
    optimization problem, then they have the same variables and the same
    objective function ($\psucc$), with the Optimal UQSD problem posing extra
    constraints on the error terms ($\perr = 0$).
    This means that the optimal solution of OptUQSD is always a feasible point
    of MED+, so $(\psucc)\supmedplus \geq (\psucc)\supuqsd$.
    Then $(\psucc)\supmed = (\psucc)\supmedplus \geq (\psucc)\supuqsd$.
\end{proof}

\end{document} 